% !TEX TS-program = pdflatex
% !TEX encoding = UTF-8 Unicode

% This file is a template using the "beamer" package to create slides for a talk or presentation
% - Giving a talk on some subject.
% - The talk is between 15min and 45min long.
% - Style is ornate.

% MODIFIED by Jonathan Kew, 2008-07-06
% The header comments and encoding in this file were modified for inclusion with TeXworks.
% The content is otherwise unchanged from the original distributed with the beamer package.

\documentclass{beamer}
\setbeamercovered{invisible}



% Copyright 2004 by Till Tantau <tantau@users.sourceforge.net>.
%
% In principle, this file can be redistributed and/or modified under
% the terms of the GNU Public License, version 2.
%
% However, this file is supposed to be a template to be modified
% for your own needs. For this reason, if you use this file as a
% template and not specifically distribute it as part of a another
% package/program, I grant the extra permission to freely copy and
% modify this file as you see fit and even to delete this copyright
% notice. 


\mode<presentation>
{
  \usetheme{Warsaw}
  % or ...

  % or whatever (possibly just delete it)
}


\usepackage[english]{babel}
% or whatever

\usepackage[utf8]{inputenc}
% or whatever

\usepackage{times}
\usepackage[T1]{fontenc}
% Or whatever. Note that the encoding and the font should match. If T1
% does not look nice, try deleting the line with the fontenc.

%%% MATH RELATED

\input{macros.tex}

\title{MATH 350-2 Advanced Calculus} 
\subtitle
{} % (optional)

\author[W.R. Casper] % (optional, use only with lots of authors)
{W.R. Casper}
% - Use the \inst{?} command only if the authors have different
%   affiliation.

\institute[California State University Fullerton] % (optional, but mostly needed)
{
  Department of Mathematics\\
  California State University Fullerton}
% - Use the \inst command only if there are several affiliations.
% - Keep it simple, no one is interested in your street address.

\subject{Talks}
% This is only inserted into the PDF information catalog. Can be left
% out. 



% If you have a file called "university-logo-filename.xxx", where xxx
% is a graphic format that can be processed by latex or pdflatex,
% resp., then you can add a logo as follows:

% \pgfdeclareimage[height=0.5cm]{university-logo}{university-logo-filename}
% \logo{\pgfuseimage{university-logo}}



% Delete this, if you do not want the table of contents to pop up at
% the beginning of each subsection:
\AtBeginSubsection[]
{
  \begin{frame}<beamer>{Outline}
    \tableofcontents[currentsection,currentsubsection]
  \end{frame}
}


% If you wish to uncover everything in a step-wise fashion, uncomment
% the following command: 

%\beamerdefaultoverlayspecification{<+->}


\begin{document}

\begin{frame}
  \titlepage
\end{frame}

\begin{frame}{Outline}
  \tableofcontents
  % You might wish to add the option [pausesections]
\end{frame}

% Since this a solution template for a generic talk, very little can
% be said about how it should be structured. However, the talk length
% of between 15min and 45min and the theme suggest that you stick to
% the following rules:  

% - Exactly two or three sections (other than the summary).
% - At *most* three subsections per section.
% - Talk about 30s to 2min per frame. So there should be between about
%   15 and 30 frames, all told.

\section{Real Analysis Lecture 17}

\subsection{Continuity and Algebra}

\begin{frame}{Warm-up}
Prove that the function $f: \mathbb R^2\rightarrow\mathbb R$ defined by
$$f(x,y) = x+y$$
is continuous. (Everything Euclidean)
\end{frame}

\begin{frame}{Warm-up}
Prove that the function $f: \mathbb R^2\rightarrow\mathbb R$ defined by
$$f(x,y) = xy$$
is continuous. (Everything Euclidean)
\end{frame}

\begin{frame}{Warm-up}
Prove that the function $f: \mathbb R \times (\mathbb R\backslash\{0\})\rightarrow\mathbb R$ defined by
$$f(x,y) = x/y$$
is continuous. (Everything Euclidean)
\end{frame}

\begin{frame}{Algebra and continuity}
\begin{thm}
Let $(S,d)$ be a metric space and $f: S\rightarrow \mathbb{R}$ and $g: S\rightarrow\mathbb{R}$ be functions which are continuous at $a\in S$.
\begin{itemize}
\item $f+g$ is continuous at $a$
\item $fg$ is continuous at $a$
\item $f/g$ is continuous at $a$ as long as $g(a)\neq 0$
\end{itemize}
\end{thm}
\end{frame}

\begin{frame}{Polynomials and Rational Functions}
\begin{itemize}
\item constant functions are continuous
\item $f(x) = x$ is continuous
\item monomials $x^n$ are continuous
\item polynomials $a_0 + a_1x + \dots + a_nx^n$ are continuous
\item rational functions $p(x)/q(x)$ are continuous on their domains
\end{itemize}
\end{frame}

\subsection{Connectedness}

\begin{frame}{Connected spaces}
\begin{defn}
A metric space $(S,d)$ is called \textbf{disconnected} if there exists two nonempty, disjoint open subsets $U,V\subseteq S$ with $S=U\cup V$.\\
\pause
We call $S$ \textbf{connected} if it is not disconnected.\\
\pause
A subset $A\subseteq S$ is connected if the corresponding subspace $(A,d)$ is connected.
\end{defn}
\pause
Examples:
\begin{itemize}
\pause
\item the real line with the Euclidean metric $(\mathbb{R},d_{eucl})$ is connected
\pause
\item the real line with the discrete metric $(\mathbb{R},d_{disc})$ is disconnected
\pause
\item euclidean space $\mathbb{R}^n$ is connected
\pause
\item an interval $I\subseteq \mathbb{R}$ with the Euclidean metric is connected
\end{itemize}
\end{frame}

\begin{frame}{Challenge!}
\begin{prob}
Prove that the set $\mathbb Q\subseteq \mathbb{R}$ is disconnected.
\end{prob}
\end{frame}

\begin{frame}{Real line is connected}
\begin{thm}
The real line $\mathbb{R}$ is connected.
\end{thm}
\pause
\begin{proof}
\pause
Suppose that $\mathbb{R}$ is disconnected.\\
\pause
Then $\mathbb{R} = U\cup V$ where:
\begin{itemize}
\pause
\item $U\neq\varnothing$, $V\neq\varnothing$
\pause
\item $U\cap V = \varnothing$
\pause
\item $U$, $V$ are open
\pause
\item $U\cup V = \mathbb{R}$
\end{itemize}
\pause
In particular, $U$ is a nonempty, proper subset of $\mathbb{R}$ which is both open and closed.
\end{proof}
\end{frame}

\begin{frame}{Real line is connected}
\begin{thm}
The real line $\mathbb{R}$ is connected.
\end{thm}
\pause
\begin{proof}
\pause
By the Representation Theorem for Open Sets:
\pause
$$U = (a_1,b_1)\cup (a_2,b_2)\cup\dots $$
\pause
is a union of disjoint open intervals.\\
\pause
In particular, $b_1$ is an accumulation point of $U$ but not an element of $U$.\\
\pause
This contradicts the fact that $U$ is closed.
\end{proof}
\end{frame}

\begin{frame}{Images of connected sets}
\begin{thm}
Let $(S,d_S)$ and $(T,d_T)$ be metric spaces and $f: S\rightarrow T$ be a function.\\
If $f$ is continuous and $S$ is connected, then $f(S)$ is connected.
\end{thm}
\pause
\begin{proof}
\pause
Suppose $f(S)$ is disconnected.\\
\pause
Then $f(S) = U\cup V$ for nonempty, disjoint, open sets $U,V\subseteq f(S)$.\\
\pause
Then $f^{-1}(U)$ and $f^{-1}(V)$ are non-empty, disjoint, open sets and
$$f^{-1}(U)\cup f^{-1}(V) = f^{-1}(U\cup V) = f^{-1}(f(S)) = S.$$
\pause
This is a contadiction.
\end{proof}
\end{frame}

\begin{frame}{Any interval is connected}
\begin{thm}
Any open interval is connected
\end{thm}
\pause
\begin{proof}
\pause
We can make any interval as a continuous image:
$$f:\mathbb{R}\rightarrow\mathbb{R}$$
$$f(x) = \frac{b+a}{2} + \frac{(b-a)x}{1+x^2}$$
\pause
$$f(\mathbb{R}) = (a,b),\quad a < b$$
\end{proof}
\end{frame}

\begin{frame}{Any interval is connected}
\begin{thm}
Any closed interval is connected
\end{thm}
\pause
\begin{proof}
We can make any interval as a continuous image:
$$f:\mathbb{R}\rightarrow\mathbb{R}$$
$$f(x) = a + \frac{2(b-a)x^2}{1+x^4}$$
\pause
$$f(\mathbb{R}) = [a,b],\quad a < b$$
\end{proof}
\end{frame}

\begin{frame}{Any interval is connected}
\begin{thm}
Any half-open interval is connected
\end{thm}
\pause
\begin{proof}
\pause
We can make any interval as a continuous image:
$$f:\mathbb{R}\rightarrow\mathbb{R}$$
$$f(x) = a + \frac{(b-a)x^2}{1+x^2}$$
\pause
$$f(\mathbb{R}) = [a,b),\quad a < b$$
\pause
$$f(\mathbb{R}) = (b,a],\quad a > b$$
\end{proof}
\end{frame}

\begin{frame}{Intermediate value theorem}
\begin{thm}[Intermediate Value Theorem]
Suppose $f: [a,b]\rightarrow\mathbb{R}$ is a continuous function.\\
Then for any value $y$ between $f(a)$ and $f(b)$, there exists $c\in [a,b]$ with $f(c) =y$.
\end{thm}
\end{frame}

\begin{frame}{Intermediate value theorem}
\begin{proof}
\pause
Let $J = f([a,b])$.\\
\pause
Then $J$ is connected!\\
\pause
Let $y$ be between $f(a)$ and $f(b)$, and suppose that $y\notin J$.\\
\pause
Without loss of generality, assume $f(a) < f(b)$.\\
\pause
Then $J = U\cup V$ for
$$U = (-\infty,y)\cap J\quad\text{and}\quad V = (y,\infty)\cap J$$
\pause
These are nonempty, disjoint, and open in the subspace $J$ of $\mathbb{R}$.\\
\pause
This is a contradiction!
\end{proof}
\end{frame}

\subsection{Uniform continuity}

\begin{frame}{Uniform continuity}
\begin{defn}
Let $(S,d_S)$ and $(T,d_T)$ be metric spaces and $A\subseteq S$.
We say that a function $f: A\rightarrow T$ is \textbf{uniformly continuous} on $A$ if for all $\epsilon > 0$ there exists $\delta > 0$ such that
$$d_S(x,y)< \delta \Rightarrow d_T(f(x),f(y))< \epsilon$$
for all $x,y\in A$.
\end{defn}
\begin{itemize}
\item $f$ continuous on $A$ means $f$ is continuous at every $a\in A$
\item $\delta$ can depend on $\epsilon$ and $a$ (not on $x$)
\item with uniform continuity, $\delta$ can't depend on $a$ ...
\end{itemize}
\end{frame}

\begin{frame}{Examples}
\begin{itemize}
\item $x$ is uniformly continuous on $\mathbb R$
\item $x^2$ is NOT uniformly continuous on $\mathbb R$
\item $1/x$ is continuous on $(0,1)$
\item $1/x$ is NOT uniformly continuous on $(0,1)$
\end{itemize}
\end{frame}

\begin{frame}{Relationship with compactness}

\begin{thm}[Heine]
Let $(S,d_S)$ and $(T,d_T)$ be metric spaces and $A\subseteq S$ be compact.
Suppose that $f: A\rightarrow T$ is continuous on $A$, then $f$ is uniformly continuous on $A$.
\end{thm}
\end{frame}

\end{document}


